% Generated draft: requires mathematical and source review.
\documentclass[11pt]{article}
\usepackage[margin=25mm]{geometry}
\usepackage{fontspec}
\setmainfont{Latin Modern Roman}
\usepackage{amsmath,amssymb,mathtools,unicode-math}
\setmathfont{Latin Modern Math}
\usepackage{enumitem,needspace,xurl,hyperref}
\setlength{\parindent}{0pt}
\setlength{\parskip}{5pt}
\setlength{\emergencystretch}{3em}
\newcommand{\N}{\mathbb{N}}
\newcommand{\Z}{\mathbb{Z}}
\newcommand{\Q}{\mathbb{Q}}
\newcommand{\R}{\mathbb{R}}
\newcommand{\C}{\mathbb{C}}
\newcommand{\F}{\mathbb{F}}
\newcommand{\A}{\mathbb{A}}
\newcommand{\PP}{\mathbb{P}}
\newcommand{\E}{\mathbb{E}}
\newcommand{\eps}{\varepsilon}
\newcommand{\dd}{\,\mathrm{d}}
\begin{document}
\section*{UnsolvedMath ID: 30004275}
\textbf{Source problem code:} catalogue entry
\subsection{Definitions and mathematical statement}
In the usual unweighted interpretation, a metric graph \(G\) is obtained by identifying endpoints of edge intervals \(I_e\), indexed by an infinite edge set \(E\). Each interval carries its length measure, and the associated Hilbert space is
\[
\mathcal H_G=\bigoplus_{e\in E}L^2(I_e).
\]
This interpretation is stated explicitly as a convention: the supplied record does not specify the admissible edge intervals, countability, local finiteness, weights, or restrictions on edge lengths.

The Kirchhoff Laplacian has the edgewise differential expression
\[
(Hf)_e=-f_e''.
\]
At a vertex \(v\) of finite degree, the Kirchhoff conditions require all incident endpoint values to agree and require
\[
\sum_{e\sim v}\partial_{\nu_e}f_e(v)=0,
\]
where \(\partial_{\nu_e}\) denotes differentiation from the vertex into the incident edge, and incidences are counted separately, including the two incidences of a loop. At a vertex of degree one this condition is the Neumann condition. The operator domain must specify suitable edgewise Sobolev regularity and global integrability in addition to these vertex conditions. Conditions at vertices of infinite degree, if allowed, require further specifications absent from the record.

For a specified densely defined symmetric realization \(H\) in \(\mathcal H_G\), define its deficiency indices by
\[
n_+(H)=\dim\ker(H^*-iI),\qquad
n_-(H)=\dim\ker(H^*+iI).
\]
Self-adjointness means equality of operators \(H=H^*\), including their domains. Essential self-adjointness instead means \(\overline H=H^*\). A self-adjoint extension is an operator \(A\) on \(\mathcal H_G\) satisfying
\[
H\subseteq A=A^*\subseteq H^*.
\]

The requested research problem is to find geometric or combinatorial hypotheses on infinite metric graphs that guarantee self-adjointness of their intended Kirchhoff Laplacian \(H\), and, for graphs whose intended realization is not self-adjoint, to determine \(n_+(H)\), \(n_-(H)\), and all its self-adjoint extensions. The record requests sufficient criteria; it does not explicitly demand a necessary-and-sufficient characterization.

This formulation remains conditional on identifying the intended graph class and operator domain. In particular, the source does not distinguish a compactly supported initial Kirchhoff operator, its closure, and a maximal realization, nor does it specify whether the intended target is self-adjointness or essential self-adjointness of an initial operator. These choices cannot be silently identified and must be resolved before this becomes a single precise operator-theoretic problem.
\subsection{Short English statement}
For an infinite metric graph, seek geometric or combinatorial conditions ensuring that its Kirchhoff Laplacian is self-adjoint. When it is not, determine its deficiency indices and self-adjoint extensions. The supplied record leaves the graph class and operator domain unspecified, so the precise target requires clarification during mathematical review.
\subsection{Sources}
\begin{itemize}
\item Local UnsolvedMath registry: permanent ID \(30004275\), source problem code \texttt{catalogue entry}, title \emph{Self-Adjointness and Extensions of Infinite Metric-Graph Laplacians}, collection \emph{Oberwolfach Reports Open Problems}. No individual proposer is supplied.
\item Supplied report citation: \emph{Mini-Workshop: Self-adjoint Extensions in New Settings} (2020), Oberwolfach Reports, pp.\ 2911--2950, DOI \texttt{10.4171/owr/2019/47}. \url{https://doi.org/10.4171/owr/2019/47}. The background attributes the problem to workshop open-problem discussions documented in Oberwolfach Reports. The record separately supplies proposed year 2019 and the tag \texttt{owr-year:2019}; the relationship between those dates and the citation year is unresolved.
\item Supplied literature reference, labeled ``Nicolussi--Kostenko--Mugnolo, self-adjoint and Markovian extensions of infinite quantum graphs'': \url{https://arxiv.org/abs/1911.04735}. Author attribution and descriptive wording are retained as supplied, without independent bibliographic verification.
\item Supplied literature reference, labeled ``original Oberwolfach report'': \url{https://doi.org/10.4171/owr/2019/47}.
\item Literature assessment dated exactly \texttt{2026-08-22}: the supplied curation labels the problem open and states that completeness, bounded degree, intrinsic metrics, and graph-end methods provide important sufficient criteria and describe extensions, while a general self-adjointness criterion for arbitrary infinite metric graphs remains open. This is supplied dated curation, not an independently verified assessment. The background expressly requests independent verification of time-sensitive status and cited claims.
\item The supplied registry labels the statement \texttt{corrected\_verified}, explaining that mathematical content was preserved while notation, typography, duplicate clauses, or omitted definitions were repaired. This is a registry assertion and does not constitute mathematical verification of this draft.
\end{itemize}
\end{document}
